\chapter*{Contribuciones Personales}
\label{cap:contribucionesPersonales}
\addcontentsline{toc}{chapter}{Contribuciones Personales}
\markboth{Contribuciones Personales}{Contribuciones Personales}

El presente Trabajo de Fin de Grado ha sido desarrollado en pareja durante el curso 2025-2026
por Juan Andrés Hibjan Cardona y Leonardo Prado de Souza. La división del trabajo se ha
articulado en torno a la especialización de cada autor, manteniendo un reparto global
aproximadamente equilibrado: la integración con fuentes externas de datos y la infraestructura
de la capa de servicios han recaído principalmente en Juan Andrés, mientras que el desarrollo
del cliente web, el diseño de la experiencia de usuario y la evaluación empírica con usuarios
han recaído principalmente en Leonardo. La lógica nuclear del motor de navegación,
especialmente la máquina de estados desarrollada durante la fase de prototipo, se ha trabajado
de manera conjunta en sesiones de diseño compartidas y posteriormente implementada y revisada
al cincuenta por ciento. Ambos hemos participado en igual medida en las decisiones
arquitectónicas, en la redacción cruzada de los capítulos y en la revisión final de la memoria.

\section*{Juan Andrés Hibjan Cardona}

Mi aportación al proyecto se ha concentrado en torno a tres ejes complementarios: la
\textbf{conexión con fuentes de datos externas}, la \textbf{infraestructura de la capa de
servicios} del motor de navegación y la \textbf{documentación técnica operativa} en los
apéndices de la memoria. La distribución del trabajo con mi compañero ha sido equilibrada en
proporción global pero diferenciada en especialización, lo que ha permitido avanzar en paralelo
durante la mayor parte del proyecto.

En la fase inicial, dedicada al prototipo en Java, mi responsabilidad ha sido la \textbf{capa
de integración de datos}. He diseñado y escrito los \textit{scripts} de ingesta para los dos
dominios sobre los que se ha validado el motor. Para TMDB, los \textit{scripts} de Python
(\texttt{scripts/TMDB/pull.py} y \texttt{scripts/TMDB/process.py}) consultan la API pública de
\textit{The Movie Database}, transforman las respuestas según el formato declarado en
\texttt{format.json} y producen el volcado intermedio en JSON Lines que sirve de fuente única
de verdad para la carga posterior. Para DBLP, el procesador escrito en Go
(\texttt{scripts/DBLP/main.go}) aplica las heurísticas de filtrado por institución sobre el
volcado público completo y produce el \textit{dataset} reducido que se carga finalmente en la base de
datos. He desarrollado además el \textit{script} unificado de poblado
(\texttt{scripts/populate\_db.py} y su variante \texttt{populate\_db\_jsonl.py}), que lee los
volcados intermedios y los inserta en el esquema relacional definido en
\texttt{database/01\_schema.sql}.

En el \textit{backend} Java, mi parte ha sido la \textbf{capa de persistencia y los controladores
infraestructurales}. Concretamente, he sido responsable de \texttt{Database\allowbreak Service}
y de la lógica de conexión a PostgreSQL, del controlador de sesión \texttt{Session\allowbreak Servlet},
del controlador central de transiciones \texttt{Navigation\allowbreak Servlet} (que orquesta todas
las acciones del motor) y del \texttt{CORSFilter} que permite el consumo desde el cliente web
en el escenario de desarrollo. También he implementado los servlets de cabecera del dominio
(\texttt{DatasetsServlet}, \texttt{CollectionsServlet}, \texttt{EntityServlet},
\texttt{EntitiesServlet} y \texttt{ResourceServlet}) que sirven el contexto sobre el que opera
el resto del motor.

En la lógica de la \textbf{máquina de estados} (\texttt{State}, \texttt{StateManager} y
\texttt{Link}) el trabajo ha sido conjunto. Hemos diseñado la representación del estado, las
transiciones, el historial y los invariantes en sesiones de trabajo compartidas, repartiéndonos
después las implementaciones concretas y revisándolas entre ambos. Estimo que el reparto de
líneas escritas en esa parte está aproximadamente al cincuenta por ciento.

En el \textit{frontend}, mi aportación ha sido menor y se ha limitado a colaborar en la primera versión
de la capa de comunicación con la API (\texttt{src/api.js}) y en la integración con los servlets
durante la depuración de los problemas de CORS y de manejo de sesión. El resto del \textit{frontend}
(componentes, diseño visual, evaluación con usuarios) ha sido responsabilidad de mi compañero.

Para el despliegue, he sido responsable de los \texttt{Dockerfile} del \textit{frontend} y del \textit{backend},
del \texttt{docker-compose.dev.yml} y \texttt{docker-compose.prod.yml}, de la configuración de
variables de entorno (\texttt{.env.db}, \texttt{.env.thirdparty}) y del \texttt{nginx.conf} del
\textit{frontend} en producción.

En la redacción de la memoria, he liderado el \textbf{capítulo de Arquitectura y Metodología},
la sección de \textbf{Modelo de datos} dentro del capítulo de Desarrollo, los dos
\textbf{apéndices técnicos} (Integración e ingesta de fuentes de datos, y Despliegue del
sistema), y he revisado y completado el capítulo de Estado de la Cuestión. He participado
también en la depuración cruzada del resto del documento y en las decisiones editoriales
conjuntas, en particular en la selección de citas y la organización de la bibliografía.

\section*{Leonardo Prado de Souza}

Mi aportación al proyecto se ha concentrado en torno a tres ejes complementarios: el
\textbf{desarrollo completo del \textit{frontend}}, el \textbf{diseño de la experiencia de usuario} y
la \textbf{evaluación empírica con usuarios externos}. La distribución del trabajo con mi
compañero ha sido equilibrada en proporción global pero diferenciada en especialización, lo que
ha permitido avanzar en paralelo durante la mayor parte del proyecto.

El \textbf{\textit{frontend}} ha sido íntegramente mi responsabilidad. He diseñado e implementado la
totalidad del cliente web en JavaScript modular, sin dependencias de \textit{framework}, sobre
una arquitectura orientada a componentes. La organización en componentes
(\texttt{src/components/Combobox.js}, \texttt{Modal.js}, \texttt{PaginationHandler.js},
\texttt{FilterTags.js}, \texttt{PanelDragOrder.js}, \texttt{Skeleton.js}, \texttt{Toast.js}),
la separación de estilos (\texttt{src/styles/} con subdivisión por componentes y por capa
estructural), la lógica de estado en \texttt{main.js} y la integración con la API mediante
\texttt{api.js} son resultado de mi trabajo. El sistema de filtros con \textit{combobox}
buscables, la diferenciación visual de filtros include/exclude con activación inline en el
propio \textit{chip}, el grafo horizontal de historial de navegación, el sistema de
\textit{drag and drop} para reordenar paneles con persistencia en \textit{localStorage} y
\textit{auto-scroll}, las notificaciones \textit{toast} y los esqueletos de carga son
componentes que he diseñado e implementado de manera independiente.

El \textbf{diseño visual} (paleta de colores, tipografía, jerarquía de paneles, sistema de
\textit{breadcrumb}, diferenciación entre acciones recuperables y destructivas en la barra
superior, estados de \textit{hover}/\textit{focus} y \textit{feedback} visual de las interacciones) es
también obra mía. La paleta y la organización final son resultado de varias iteraciones
documentadas, parte de ellas guiadas por los hallazgos del estudio de usabilidad descrito más
adelante.

He liderado la \textbf{evaluación de usabilidad}, desde el diseño del protocolo hasta la
conducción de las sesiones y la codificación posterior. He elaborado el guion completo de
evaluación (reproducido en el apéndice~\ref{Appendix:Guion}), gestionado el reclutamiento de
los participantes, conducido las cinco sesiones presenciales, transcrito las observaciones y
citas literales, calculado las métricas cuantitativas (SUS, SEQ, tasa de éxito, tiempo por
tarea) y derivado el catálogo de problemas detectados con la escala de gravedad de Nielsen. Las
iteraciones de interfaz que se derivaron de la evaluación las hemos implementado conjuntamente,
pero el diseño y la conducción del estudio han sido mi responsabilidad exclusiva.

En el \textit{backend} Java, durante la fase del prototipo, mi parte ha sido la lógica de
\textbf{filtros facetados y operaciones de conjunto}. He sido responsable de los servlets
\texttt{MetadataFacetsServlet}, \texttt{ReferenceFacetsServlet}, \texttt{LinkFacetsServlet} y
\texttt{UnionServlet}, así como de la parte de \texttt{NavigationDAO} correspondiente a la
traducción de las primitivas de filtrado a consultas SQL parametrizadas y a la construcción
incremental de cláusulas \texttt{WHERE} a partir del estado.

En la lógica central de la \textbf{máquina de estados} (\texttt{State}, \texttt{StateManager}
y \texttt{Link}) el trabajo ha sido conjunto con mi compañero, en sesiones compartidas de
diseño seguidas de implementación repartida y revisión cruzada. Estimo que el reparto de líneas
escritas en esa parte está aproximadamente al cincuenta por ciento.

En la redacción de la memoria, he liderado el \textbf{capítulo de Pruebas y Validación}, el
\textbf{apéndice del guion de evaluación de usabilidad}, las secciones de \textbf{Desarrollo
del \textit{frontend}} y \textbf{Operaciones de conjuntos} dentro del capítulo de Desarrollo, y he
participado en la redacción del capítulo de Introducción y en las líneas de trabajo futuro.
He coordinado además las revisiones tipográficas y de estilo del documento global, así como la
generación de las figuras vectoriales con PGFplots para los resultados cuantitativos.
